% Ch8 WS3 -- Lewis structures, formal charge, resonance, VSEPR + FRQ.
\documentclass{apchem-worksheet}

\apcunitword{Chapter}
\apcunit{8}
\apcunittitle{Bonding: General Concepts}
\apcdoctitle{Worksheet 3 \textbullet\ Structures, Resonance, VSEPR \& FRQ}
\apcsubtitle{Zumdahl \S8.9--8.13 \textbullet\ electron count first, always}

\begin{document}
\wsheader[Write the total valence electron count before drawing anything ---
an uncounted structure cannot be checked. Every geometry answer needs its
domain count shown.]

\problem[]{Electron counts and Lewis structures:}
\begin{parts}
  \item \ce{NH4+} \quad count: \answerblank[1.8cm]{8}
        \workspace[2cm]
        \keyonly{\begin{apcanswer}N central, four N--H single bonds, no lone
        pairs, brackets with $+$.\end{apcanswer}}
  \item \ce{SO3} \quad count: \answerblank[1.8cm]{24}
        \workspace[2cm]
        \keyonly{\begin{apcanswer}S central, three resonance forms with one
        S=O and two S--O (or all three drawn double); trigonal
        planar.\end{apcanswer}}
  \item \ce{XeF4} \quad count: \answerblank[1.8cm]{36}
        \workspace[2cm]
        \keyonly{\begin{apcanswer}Xe central with four Xe--F bonds and TWO
        lone pairs on Xe --- expanded octet, 12 electrons around
        Xe.\end{apcanswer}}
\end{parts}

\problem[]{Octet exceptions. Name the class and explain:}
\begin{parts}
  \item \ce{NO2}: \answerblank[6cm]{odd-electron species (17 e$^-$)}
  \item \ce{BeCl2}: \answerblank[6.4cm]{fewer than an octet (Be gets 4)}
  \item \ce{PCl5}: \answerblank[6cm]{expanded octet (period 3)}
  \item Explain why \ce{NCl5} does not exist.
        \answerlines[2]{Nitrogen is in period 2 and cannot exceed eight
        electrons, so it cannot form five bonds.}
\end{parts}

\problem[]{Formal charge:}
\begin{parts}
  \item Compute FC on each atom in \ce{NO3-} for a form with one N=O and
        two N--O. \answerlines[2]{N: $5 - 0 - 4 = +1$. Double-bonded O:
        $6-4-2 = 0$. Each single-bonded O: $6-6-1 = -1$. Sum
        $= -1$ \checkmark}
  \item State the two selection rules for choosing among candidate
        structures. \answerlines[2]{Formal charges as close to zero as
        possible; any negative formal charge placed on the more
        electronegative atom.}
\end{parts}

\problem[]{Resonance:}
\begin{parts}
  \item Draw all resonance forms of \ce{CO3^2-} (24 e$^-$).
        \workspace[3cm]
        \keyonly{\begin{apcanswer}Three forms with the C=O rotating among
        the three oxygens; brackets with $2-$.\end{apcanswer}}
  \item Bond order of each C--O bond: \answerblank[2cm]{$4/3$}
  \item State the experimental evidence that carbonate is a hybrid rather
        than a molecule alternating between forms.
        \answerlines[2]{All three C--O bonds are measured to be identical in
        length and intermediate between a single and a double bond ---
        alternation would give two distinct lengths.}
\end{parts}

\problem[]{Complete the VSEPR chain for each species:}
\begin{parts}
  \item \ce{PCl3}: domains \answerblank[2.0cm]{4 (1 LP)}, geometry
        \answerblank[3.9cm]{trigonal pyramidal}, angle
        \answerblank[2.2cm]{$\approx107^\circ$}
  \item \ce{CS2}: domains \answerblank[1.6cm]{2}, geometry
        \answerblank[3.4cm]{linear}, angle \answerblank[2.2cm]{180$^\circ$}
  \item \ce{H2S}: domains \answerblank[2.0cm]{4 (2 LP)}, geometry
        \answerblank[3.4cm]{bent}, angle
        \answerblank[2.4cm]{$<109.5^\circ$}
  \item \ce{SO3}: domains \answerblank[1.6cm]{3}, geometry
        \answerblank[3.4cm]{trigonal planar}, angle
        \answerblank[2.2cm]{120$^\circ$}
\end{parts}

\problem[]{\textbf{FRQ (10 points).} Sulfur and carbon each form a
triatomic oxide: \ce{SO2} and \ce{CO2}.}
\begin{parts}
  \item Give the total valence electron count for each.
        \answerlines[1]{\ce{SO2}: 18. \ce{CO2}: 16.}
  \item Draw a complete Lewis structure for each, showing all lone pairs.
        \workspace[3.2cm]
        \keyonly{\begin{apcanswer}\ce{CO2}: O=C=O, two lone pairs per O,
        none on C. \ce{SO2}: bent, one S=O and one S--O (resonance), two
        lone pairs on each terminal O and ONE LONE PAIR ON
        S.\end{apcanswer}}
  \item State the electron-domain count, molecular geometry, and
        approximate bond angle for each.
        \answerlines[3]{\ce{CO2}: 2 domains, linear, 180$^\circ$.
        \ce{SO2}: 3 domains, bent, slightly less than 120$^\circ$.}
  \item One is polar and one is not. Identify which and justify using bond
        dipoles and geometry. \answerlines[2]{\ce{SO2} is polar --- its bent
        geometry prevents the two S--O dipoles from cancelling. \ce{CO2} is
        nonpolar --- its linear geometry cancels them exactly.}
  \item The two S--O bonds in \ce{SO2} are found to be equal in length.
        Explain. \answerlines[2]{\ce{SO2} is a resonance hybrid of two
        equivalent forms, so each bond has the same intermediate bond order
        and therefore the same length.}
\end{parts}

\begin{apcnote}[Rubric --- score yourself, 10 points]
\textbf{(a) 1 pt} both counts \textbullet\ \textbf{(b) 2 pts} 1 \ce{CO2}
correct; 1 \ce{SO2} correct \emph{including the lone pair on S} --- omitting
it earns 0 for this point \textbullet\ \textbf{(c) 3 pts} 1 both domain
counts, 1 both geometries, 1 both angles \textbullet\ \textbf{(d) 2 pts}
1 identifies \ce{SO2}; 1 argues cancellation in both directions
\textbullet\ \textbf{(e) 2 pts} 1 names resonance/hybrid, 1 links equal
bond order to equal length --- ``the bond alternates'' earns 0.
\end{apcnote}

\begin{spiralstrip}{Chapter 8 \textbullet\ blocks 1--3}
  \item $\Delta H$ from bond energies:
        \answerblank[4cm]{broken $-$ formed}
  \item Bond order of \ce{O3} O--O bonds: \answerblank[1.8cm]{$3/2$}
  \item Higher lattice energy: \ce{KBr} or \ce{CaO}?
        \answerblank[2.4cm]{\ce{CaO}}
  \item Electrons in \ce{ClO3-}: \answerblank[2cm]{26}
  \item Can period-2 atoms expand their octet? \answerblank[1.6cm]{no}
\end{spiralstrip}

\end{document}
